\documentclass[journal,onecolumn]{IEEEtran}

% ── Encabezado ──────────────────────────────────────────────
\usepackage{fancyhdr}
\pagestyle{fancy}
\fancyhf{}
\renewcommand{\headrulewidth}{0.4pt}
\fancyhead[L]{\small Fundamento Teórico — Bloque FT.2}
\fancyfoot[C]{\thepage}

% ============================================================
% IDIOMA Y CODIFICACIÓN  (mover ANTES de csquotes)
% ============================================================
\usepackage[utf8]{inputenc}
\usepackage[T1]{fontenc}
\usepackage[spanish,es-tabla]{babel}

% ============================================================
% CITAS Y TEXTO
% ============================================================
\usepackage{csquotes}

% ============================================================
% ESPACIADO Y LISTAS
% ============================================================
\usepackage{parskip}
\setlength{\parskip}{4pt}
\usepackage{enumitem}

% ============================================================
% IDIOMA Y CODIFICACIÓN
% ============================================================
\usepackage[spanish,es-tabla]{babel}
\usepackage[utf8]{inputenc}
\usepackage[T1]{fontenc}

% ============================================================
% PAPEL Y GEOMETRÍA
% ============================================================
\usepackage{geometry}
\geometry{a4paper, top=2.5cm, bottom=2.5cm, left=2.5cm, right=2.5cm}

% ============================================================
% IMÁGENES
% ============================================================
\usepackage{graphicx}
\graphicspath{{Figuras/}}

% ============================================================
% TABLAS
% ============================================================
\usepackage{booktabs}
\usepackage{tabularx}
\usepackage{longtable}

% ============================================================
% MATEMÁTICAS
% ============================================================
\usepackage{amsmath}
\usepackage{amssymb}

% ============================================================
% CÓDIGO FUENTE
% ============================================================
\usepackage{listings}

% ============================================================
% FORMATO ACADÉMICO
% ============================================================
\usepackage{microtype}
\usepackage{float}

% ============================================================
% CAPTIONS (TÍTULOS DE FIGURAS/TABLAS)
% ============================================================
\usepackage{caption}
\captionsetup{font=small, labelfont=bf}

% ============================================================
% REFERENCIAS / HIPERVÍNCULOS
% ============================================================
\usepackage[hidelinks]{hyperref}

% ============================================================
% NUMERACIÓN Y NOMBRES EN ESPAÑOL
% ============================================================
\renewcommand{\thetable}{\arabic{table}}
\def\tablename{Tabla}

\renewcommand{\thesection}{\arabic{section}}
\renewcommand{\thesubsection}{\thesection.\arabic{subsection}}
\renewcommand{\thesubsubsection}{\thesubsection.\arabic{subsubsection}}

% ============================================================
% ESTILO DE TÍTULOS DE SECCIÓN
% ============================================================
\usepackage{titlesec}
\titleformat{\section}{\normalfont\Large\bfseries}{\thesection}{1em}{}
\titleformat{\subsection}{\normalfont\large\bfseries}{\thesubsection}{1em}{}
\titleformat{\subsubsection}{\normalfont\normalsize\bfseries}{\thesubsubsection}{1em}{}

% ============================================================
% CADA SECCIÓN PRINCIPAL (\section) INICIA EN PÁGINA NUEVA
% ============================================================
\usepackage{etoolbox}
\pretocmd{\section}{\clearpage}{}{}

% ============================================================
% BIBLIOGRAFÍA (BIBLATEX + BIBER)
% ============================================================
\usepackage[
backend=biber,
style=numeric,
sorting=none
]{biblatex}
\addbibresource{referencias_FT2.bib}

% ============================================================
% DOCUMENTO
% ============================================================
\begin{document}

% ============================================================
% FUNDAMENTO TEÓRICO — BLOQUE FT.2
% ============================================================

\section{Fundamentos de la especificación y el ciclo de vida}

\subsection{Ingeniería de requisitos y fundamentos de los requisitos}

La ingeniería de requisitos comprende el conjunto de actividades orientadas a identificar, analizar,
documentar, validar y gestionar las necesidades que un sistema debe satisfacer, de modo que dichas
necesidades puedan servir de base verificable para el resto de las actividades de desarrollo
\cite{ieee29148,pohl2025}. Esta disciplina no se limita a recoger enunciados de los interesados: los
organiza mediante un marco de trabajo que distingue el contexto en que opera el sistema, el propio
proceso de obtención y negociación, y los artefactos documentados que resultan de ese proceso
\cite{pohl2025}. En el caso de MundiPets, este marco es lo que permite tratar los requisitos no como
una lista informal recogida de los futuros usuarios de la plataforma, sino como un conjunto
estructurado que puede auditarse, corregirse y cerrarse con criterios objetivos.

Un requisito individual, para ser útil al proceso de desarrollo, debe reunir un conjunto de
características de calidad reconocidas: debe ser necesario, no ambiguo, completo, singular, factible,
verificable y trazable, entre otras \cite{ieee29148}. Estas características no son exigencias
puramente formales: son la condición que hace posible que, más adelante, un requisito pueda auditarse
con una métrica, enlazarse a un caso de uso y a un caso de prueba, y defenderse ante un tribunal con
un artefacto concreto que lo respalde. Por ello, el fundamento de la especificación de MundiPets no
puede separarse del fundamento de la calidad del requisito individual.

\subsection{Elicitación, análisis y validación de requisitos}

Antes de que un requisito pueda especificarse con calidad, debe elicitarse, analizarse y negociarse
con los interesados del sistema. Pohl describe estas actividades como parte de un ciclo continuo en el
que la elicitación aporta el conocimiento disperso entre los interesados, el análisis lo estructura y
depura, y la negociación resuelve los conflictos de intereses que surgen entre distintos puntos de
vista sobre el mismo sistema \cite{pohl2025}. El SWEBOK sitúa estas mismas actividades dentro del área
de conocimiento de Requisitos de Software, subrayando que la elicitación y el análisis no son un paso
previo desechable, sino una fuente de retroalimentación continua sobre la propia especificación
\cite{swebok2024}. El temario del IREB para el nivel Foundation coincide en este punto: la obtención,
la documentación y la validación de requisitos se presentan como prácticas interdependientes, no como
fases estancas de un proceso lineal \cite{ireb2026}.

La validación cierra este ciclo al comprobar que los requisitos documentados representan realmente lo
que los interesados necesitan, y no únicamente lo que fue posible transcribir en una primera
iteración \cite{ieee29148,pohl2025}. En el proyecto integrador, esta actividad se corresponde con la
inspección formal ejecutada en la PE4 y con la re-inspección de la PE5 sobre el ERS ya corregido
(Sección 5): ambas son, en términos de este marco teórico, instancias concretas de validación de
requisitos aplicadas al caso de MundiPets.

\subsection{Especificación de requisitos y estructura del ERS}

La especificación traduce el resultado de la elicitación, el análisis y la validación en una
representación documentada, estructurada y verificable de las necesidades y características que debe
satisfacer el sistema \cite{ieee29148}. La ISO/IEC/IEEE 29148:2018 define el contenido y la estructura
exigibles a una Especificación de Requisitos de Software (SRS), estableciendo las cláusulas que un
documento de este tipo debe contemplar y las características de calidad aplicables tanto al conjunto
de requisitos como a cada requisito individual \cite{ieee29148}. El SWEBOK, por su parte, ubica la
especificación como una de las prácticas centrales del área de Requisitos de Software y la vincula
explícitamente con la gestión de requisitos y con la trazabilidad, entendida esta última como la
capacidad de seguir la vida de un requisito en ambas direcciones a lo largo del documento y del resto
de artefactos del proyecto \cite{swebok2024}. Pohl aporta a esta misma idea el fundamento conceptual
de los tipos de artefactos de requisitos objetivos, escenarios y requisitos orientados a la
solución que la especificación debe documentar de forma coherente entre sí \cite{pohl2025}.

Desde esta perspectiva, el ERS de MundiPets no constituye únicamente un documento descriptivo del
sistema, sino un artefacto que fija una referencia verificable para el desarrollo posterior: cada
requisito documentado en él debe poder auditarse contra los criterios de completitud, consistencia y
verificabilidad que se aplican en la Sección 8, y debe poder trazarse hacia los casos de uso, los
modelos y los casos de prueba que lo realizan, tal como exige la cadena de trazabilidad de la
Sección 6.

\subsection{Requisitos dentro del ciclo de vida del sistema y del software}

Los requisitos no se especifican de manera aislada: se insertan dentro de un ciclo de vida que
abarca, según el nivel de abstracción, al sistema completo o específicamente al software. La
ISO/IEC/IEEE 15288:2023 define un conjunto de procesos de ciclo de vida aplicables a sistemas creados
por el ser humano y sitúa la definición de las necesidades y requisitos de las partes interesadas, así
como el análisis de los requisitos del sistema, dentro de los procesos técnicos del ciclo de vida
\cite{iso15288}. Esta norma proporciona el encuadre bajo el cual MundiPets se concibe como un sistema
con procesos de definición, desarrollo, verificación, validación y evolución propios, de los cuales la
especificación de requisitos es una entrada temprana y determinante.

Cuando el interés se concentra específicamente en el software, la ISO/IEC/IEEE 12207:2017 en su
edición de 2017, vigente al momento de redactar este informe define el marco de procesos de ciclo
de vida del software, incluidos los relacionados con su desarrollo, su mantenimiento y su operación, y
comparte con la ISO/IEC/IEEE 15288:2023 una estructura de procesos armonizada para facilitar su uso
conjunto \cite{iso12207}. Para MundiPets, esto significa que los requisitos del ERS no se cierran en la
PE5 de forma definitiva: quedan insertados en un ciclo de vida de software que contempla su evolución
futura, su mantenimiento y su eventual modificación, razón por la cual la trazabilidad y la
modificabilidad auditadas en la Sección 8 no son un ejercicio académico aislado, sino una condición
para que ese ciclo de vida posterior sea manejable.

\subsection{Prácticas reconocidas de ingeniería de requisitos}

Más allá de las normas de proceso, el equipo se apoya en dos cuerpos de conocimiento que consolidan
las prácticas reconocidas de la disciplina. El SWEBOK Guide, en su versión 4.0, dedica un área de
conocimiento completa a los Requisitos de Software y organiza sus prácticas en torno a los
fundamentos del requisito, la elicitación, el análisis, la especificación, la validación y la gestión
de requisitos, incluida la trazabilidad \cite{swebok2024}. El IREB, por su parte, sostiene el temario
CPRE Foundation Level en su versión 3.3.0, que estructura el conocimiento básico exigible a un
ingeniero de requisitos en torno a los interesados, las técnicas de elicitación, la documentación de
requisitos, su validación, su gestión y su comunicación \cite{ireb2026}. Ambos cuerpos de conocimiento
coinciden con Pohl y con la propia ISO/IEC/IEEE 29148:2018 en un punto central para MundiPets: ninguna
de estas prácticas se sostiene si el requisito documentado no puede verificarse objetivamente, lo que
en esta práctica se traduce en la exigencia de un criterio de aceptación en formato
Dado--Cuando--Entonces para cada requisito del sistema (Sección 3.2).

\subsection{Integración de los fundamentos para el cierre del ERS de MundiPets}

Las seis fuentes utilizadas en este bloque cumplen funciones complementarias y no redundantes dentro
de la argumentación que sustenta el cierre del ERS de MundiPets. La ISO/IEC/IEEE 29148:2018 aporta el
contenido y la estructura exigibles a la especificación misma: es la norma contra la cual se verifica
que el documento entregado sea, formalmente, un SRS y no una colección de notas de reunión
\cite{ieee29148}. La ISO/IEC/IEEE 15288:2023 sitúa esa especificación dentro del ciclo de vida del
sistema MundiPets en su conjunto, mientras que la ISO/IEC/IEEE 12207:2017 hace lo propio para el ciclo
de vida específico del software que se construirá a partir de ese ERS \cite{iso15288,iso12207}; ambas
normas comparten estructura de procesos, de modo que no compiten entre sí, sino que delimitan el nivel
de abstracción  (sistema o software ) desde el que se lee la misma cadena de requisitos. Pohl
aporta los fundamentos y las actividades de la ingeniería de requisitos que preceden a la
especificación formal (elicitación, análisis, negociación, validación) y que explican por qué el
contenido del ERS es defendible y no arbitrario \cite{pohl2025}. El SWEBOK aporta el cuerpo de
conocimiento y las prácticas reconocidas de la disciplina que enmarcan estas actividades dentro de la
ingeniería de software en general \cite{swebok2024}, y el temario CPRE Foundation Level del IREB aporta
las prácticas estructuradas y verificables de ingeniería de requisitos que un profesional certificable
debe poder aplicar \cite{ireb2026}.

La integración de estos seis fundamentos permite sostener que el cierre del ERS de MundiPets no es un
acto administrativo de fin de semestre, sino la culminación fundamentada de un proceso: los requisitos
fueron elicitados y analizados conforme a las prácticas descritas por Pohl, el SWEBOK y el IREB
\cite{pohl2025,swebok2024,ireb2026}; fueron especificados con la estructura y las características de
calidad que exige la ISO/IEC/IEEE 29148:2018 \cite{ieee29148}; fueron auditados, corregidos y trazados
de extremo a extremo conforme a esas mismas características (Secciones 6 y 8); y quedan insertados en
un ciclo de vida de sistema y de software reconocido por la ISO/IEC/IEEE 15288:2023 y la
ISO/IEC/IEEE 12207:2017 \cite{iso15288,iso12207}, que sostiene su evolución más allá de esta práctica.
Sobre esa base, y no sobre una declaración de intenciones, el equipo considera que el ERS final de
MundiPets reúne el fundamento teórico y normativo necesario para su cierre.

\printbibliography[heading=none]

\end{document}
